\section{Datasets presentation}

\begin{frame}
	\frametitle{Presentation of the datasets used}
	
	3 datasets:
	\begin{enumerate}
		\item Census grid 2021
		\item Geographic accessibility
		\item Communes
	\end{enumerate}
	
\end{frame}


\begin{frame}
	\frametitle{Filtering the datasets}
	
	Definition of table \texttt{it\_communes} and indexes for all the tables.
	
\end{frame}


\begin{frame}[containsverbatim]
	\frametitle{Calculation of the population by municipality}
	
	The formula used:
	\begin{equation}
		P_{C} = P_{G} \sum \frac{A(C \cap G)}{A(G)}
	\end{equation}
	
	and the query performed:
	\begin{figure}[h]
	\begin{minted}{sql}
WITH pop_calc AS (
	SELECT 
		c."COMM_ID",
		ROUND(SUM(ST_Area(ST_Intersection(g.geom, c.geom)) / 1000000 * g."T")) AS pop
	FROM census_2021 g JOIN it_communes c ON ST_Intersects(g.geom, c.geom)
	GROUP BY c."COMM_ID"
)
UPDATE it_communes c
SET pop = p.pop
FROM pop_calc p
WHERE c."COMM_ID" = p."COMM_ID";
	\end{minted}
	\end{figure}
	
\end{frame}


\begin{frame}[containsverbatim]
	\frametitle{Calculation of mean accessibility value}
	
	Computation of the mean accessibility times by commune to the closest healthcare structure and to the 3 closest ones. For years 2020 and 2023.
	\begin{figure}[h]
	\begin{minted}{sql}
WITH avg_t AS (
	SELECT
		c."COMM_ID",
		AVG(g.health_2020_n1) AS mean_health_2020_n1,
		AVG(g.health_2020_n3) AS mean_health_2020_n3,
		...
	FROM it_communes c JOIN grid_accessibility_health g ON st_intersects(c.geom, g.geom)
	GROUP BY c."COMM_ID"
)
UPDATE it_communes as c
SET
	mean_health_2020_n1 = ROUND(t.mean_health_2020_n1::NUMERIC, 2),
	...
FROM avg_t t
WHERE t."COMM_ID"=c."COMM_ID";
	\end{minted}
	\end{figure}
	
\end{frame}







